% GENERATED FILE. Edit canonical lessons, then run npm run generate:books.
% canonical-source-hash: fnv1a64:9b0f6e2310bca7eb
% canonical-lessons: MR-C348-sambandh, MR-C348-samya, MR-C348-sandarbha, MR-C348-nigdit, MR-C348-yaulat

\chapter{Links and Likenesses}
\label{ch:mr-links-and-likenesses}
\hlchaptermodality{348}

\begin{chapteropening}
\textbf{By the end of this chapter:} \emph{I can say how two things are linked or alike.}

Complete the last lesson of chapter 348: I can say how two things are linked or alike.
\end{chapteropening}

\section[\texorpdfstring{\mr{संबंध} (sambandh)}{संबंध (sambandh)}]{\mr{संबंध} (sambandh) — a connection; a relationship}

\label{lesson:MR-C348-sambandh}



\begin{quote}
Before the new one: say the Marathi for neutral, then the Marathi for against.
\end{quote}

\subsection*{You'll want to know: \mr{संबंध}}
\textbf{\mr{संबंध}} — \emph{sambandh} — \textquotedblleft{}a connection; a relationship\textquotedblright{}.

\begin{quote}
\mr{या} \mr{दोन} \mr{घटनांचा} \mr{एकमेकांशी} \mr{काही} \mr{संबंध} \mr{नाही}.
\end{quote}

\begin{quote}
\emph{yā don ghaṭnāñcā ekmekāṁśī kāhī sambandh nāhī.}
\end{quote}

\begin{quote}
These two events have no connection with each other.
\end{quote}

Masculine.

X \mr{चा} Y \mr{शी} \mr{संबंध} \mr{असणे}, X is connected with Y.

\begin{cousinweb}
\textbf{\mr{संबंध}} (\emph{sambandh}) is a tatsama, taken from Sanskrit \emph{sambandha}, \textquotedblleft{}binding together, connection\textquotedblright{} (Monier-Williams).
\end{cousinweb}

\begin{grammarlens}[title={What you've built}]
One more word for giving reasons and opinions.
\end{grammarlens}

\subsection*{Your turn}
\begin{itemize}
\raggedright
  \item \emph{Say it:} \emph{sambandh}
  \item \emph{Say it:} \emph{sambandh}, once more
  \item \emph{Say it:} \emph{yā don ghaṭnāñcā ekmekāṁśī kāhī sambandh nāhī.} — These two events have no connection with each other.
  \item \emph{Recall:} say the Marathi for neutral, then the Marathi for against, then say \emph{sambandh} again
\end{itemize}

\begin{tcolorbox}[breakable,colback=teal!4,colframe=teal!35!black,title={Before you move on}]
What does \mr{संबंध} mean? (\textquotedblleft{}A connection; a relationship\textquotedblright{}.) Say it once more.
\end{tcolorbox}
\section[\texorpdfstring{\mr{साम्य} (sāmya)}{साम्य (sāmya)}]{\mr{साम्य} (sāmya) — a similarity}

\label{lesson:MR-C348-samya}



\begin{quote}
Before the new one: say the Marathi for against, then the Marathi for a connection.
\end{quote}

\subsection*{You'll want to know: \mr{साम्य}}
\textbf{\mr{साम्य}} — \emph{sāmya} — \textquotedblleft{}a similarity\textquotedblright{}.

\begin{quote}
\mr{दोन्ही} \mr{भावांच्या} \mr{आवाजात} \mr{खूप} \mr{साम्य} \mr{आहे}.
\end{quote}

\begin{quote}
\emph{donhī bhāvāñcyā āvājāt khūp sāmya āhe.}
\end{quote}

\begin{quote}
The two brothers' voices are very alike.
\end{quote}

Neuter.

Neuter; X \mr{आणि} Y \mr{मध्ये} \mr{साम्य} \mr{असणे}, X and Y are alike.

\begin{cousinweb}
\textbf{\mr{साम्य}} (\emph{sāmya}) is a tatsama, taken from Sanskrit \emph{sāmya}, \textquotedblleft{}equality, likeness\textquotedblright{} (Monier-Williams).
\end{cousinweb}

\begin{grammarlens}[title={What you've built}]
One more word for giving reasons and opinions.
\end{grammarlens}

\subsection*{Your turn}
\begin{itemize}
\raggedright
  \item \emph{Say it:} \emph{sāmya}
  \item \emph{Say it:} \emph{sāmya}, once more
  \item \emph{Say it:} \emph{donhī bhāvāñcyā āvājāt khūp sāmya āhe.} — The two brothers' voices are very alike.
  \item \emph{Recall:} say the Marathi for against, then the Marathi for a connection, then say \emph{sāmya} again
\end{itemize}

\begin{tcolorbox}[breakable,colback=teal!4,colframe=teal!35!black,title={Before you move on}]
What does \mr{साम्य} mean? (\textquotedblleft{}A similarity\textquotedblright{}.) Say it once more.
\end{tcolorbox}
\section[\texorpdfstring{\mr{संदर्भ} (sandarbha)}{संदर्भ (sandarbha)}]{\mr{संदर्भ} (sandarbha) — a context; a reference}

\label{lesson:MR-C348-sandarbha}



\begin{quote}
Before the new one: say the Marathi for a connection, then the Marathi for a similarity.
\end{quote}

\subsection*{You'll want to know: \mr{संदर्भ}}
\textbf{\mr{संदर्भ}} — \emph{sandarbha} — \textquotedblleft{}a context; a reference\textquotedblright{}.

\begin{quote}
\mr{संदर्भ} \mr{समजल्याशिवाय} \mr{या} \mr{वाक्याचा} \mr{अर्थ} \mr{कळत} \mr{नाही}.
\end{quote}

\begin{quote}
\emph{sandarbha samjalyāśivāy yā vākyācā artha kaḷat nāhī.}
\end{quote}

\begin{quote}
Without the context, the meaning of this sentence is not clear.
\end{quote}

Masculine.

\mr{संदर्भात}, regarding: \mr{या} \mr{संदर्भात}, in this connection.

\begin{cousinweb}
\textbf{\mr{संदर्भ}} (\emph{sandarbha}) is a tatsama, taken from Sanskrit \emph{saṁdarbha}, \textquotedblleft{}stringing together, composition\textquotedblright{} (Monier-Williams).
\end{cousinweb}

\begin{grammarlens}[title={What you've built}]
One more word for giving reasons and opinions.
\end{grammarlens}

\subsection*{Your turn}
\begin{itemize}
\raggedright
  \item \emph{Say it:} \emph{sandarbha}
  \item \emph{Say it:} \emph{sandarbha}, once more
  \item \emph{Say it:} \emph{sandarbha samjalyāśivāy yā vākyācā artha kaḷat nāhī.} — Without the context, the meaning of this sentence is not clear.
  \item \emph{Recall:} say the Marathi for a connection, then the Marathi for a similarity, then say \emph{sandarbha} again
\end{itemize}

\begin{tcolorbox}[breakable,colback=teal!4,colframe=teal!35!black,title={Before you move on}]
What does \mr{संदर्भ} mean? (\textquotedblleft{}A context; a reference\textquotedblright{}.) Say it once more.
\end{tcolorbox}
\section[\texorpdfstring{\mr{निगडित} (nigḍit)}{निगडित (nigḍit)}]{\mr{निगडित} (nigḍit) — linked (with)}

\label{lesson:MR-C348-nigdit}



\begin{quote}
Before the new one: say the Marathi for a similarity, then the Marathi for a context.
\end{quote}

\subsection*{You'll want to know: \mr{निगडित}}
\textbf{\mr{निगडित}} — \emph{nigḍit} — \textquotedblleft{}linked (with)\textquotedblright{}.

\begin{quote}
\mr{आरोग्य} \mr{आणि} \mr{आहार} \mr{एकमेकांशी} \mr{निगडित} \mr{आहेत}.
\end{quote}

\begin{quote}
\emph{ārogya āṇi āhār ekmekāṁśī nigḍit āhet.}
\end{quote}

\begin{quote}
Health and diet are linked with each other.
\end{quote}

Takes -\mr{शी}: X \mr{शी} \mr{निगडित}. Invariable.

\begin{cousinweb}
\textbf{\mr{निगडित}} (\emph{nigḍit}) is a tatsama, taken from Sanskrit \emph{nigaḍita}, \textquotedblleft{}chained, fettered\textquotedblright{} (Monier-Williams). Marathi uses it for things bound up with each other.
\end{cousinweb}

\begin{grammarlens}[title={What you've built}]
One more word for giving reasons and opinions.
\end{grammarlens}

\subsection*{Your turn}
\begin{itemize}
\raggedright
  \item \emph{Say it:} \emph{nigḍit}
  \item \emph{Say it:} \emph{nigḍit}, once more
  \item \emph{Say it:} \emph{ārogya āṇi āhār ekmekāṁśī nigḍit āhet.} — Health and diet are linked with each other.
  \item \emph{Recall:} say the Marathi for a similarity, then the Marathi for a context, then say \emph{nigḍit} again
\end{itemize}

\begin{tcolorbox}[breakable,colback=teal!4,colframe=teal!35!black,title={Before you move on}]
What does \mr{निगडित} mean? (\textquotedblleft{}Linked (with)\textquotedblright{}.) Say it once more.
\end{tcolorbox}
\section[\texorpdfstring{\mr{याउलट} (yāulaṭ)}{याउलट (yāulaṭ)}]{\mr{याउलट} (yāulaṭ) — on the contrary; by contrast}

\label{lesson:MR-C348-yaulat}



\begin{quote}
Before the new one: say the Marathi for a context, then the Marathi for linked.
\end{quote}

\subsection*{You'll want to know: \mr{याउलट}}
\textbf{\mr{याउलट}} — \emph{yāulaṭ} — \textquotedblleft{}on the contrary; by contrast\textquotedblright{}.

\begin{quote}
\mr{भाऊ} \mr{शांत} \mr{आहे}, \mr{याउलट} \mr{बहीण} \mr{खूप} \mr{बोलकी} \mr{आहे}.
\end{quote}

\begin{quote}
\emph{bhāū śānt āhe, yāulaṭ bahīṇ khūp bolkī āhe.}
\end{quote}

\begin{quote}
The brother is quiet; by contrast, the sister is very talkative.
\end{quote}

\mr{या} + \mr{उलट}, this + the reverse.

\begin{cousinweb}
\textbf{\mr{याउलट}} (\emph{yāulaṭ}) is built from \mr{या}, \textquotedblleft{}this\textquotedblright{} (oblique), and \mr{उलट}, \textquotedblleft{}the reverse\textquotedblright{}: the reverse of this. \mr{उलट} goes with the verb \mr{उलटणे}; CDIAL lists Marathi \emph{ulaṭṇẽ}, \textquotedblleft{}to be upset, to turn over\textquotedblright{}, under \emph{*ullaṭati}, \textquotedblleft{}turns over\textquotedblright{} (CDIAL 2368).
\end{cousinweb}

\begin{grammarlens}[title={What you've built}]
One more word for giving reasons and opinions.
\end{grammarlens}

\subsection*{Your turn}
\begin{itemize}
\raggedright
  \item \emph{Say it:} \emph{yāulaṭ}
  \item \emph{Say it:} \emph{yāulaṭ}, once more
  \item \emph{Say it:} \emph{bhāū śānt āhe, yāulaṭ bahīṇ khūp bolkī āhe.} — The brother is quiet; by contrast, the sister is very talkative.
  \item \emph{Recall:} say the Marathi for a context, then the Marathi for linked, then say \emph{yāulaṭ} again
\end{itemize}

\begin{tcolorbox}[breakable,colback=teal!4,colframe=teal!35!black,title={Before you move on}]
What does \mr{याउलट} mean? (\textquotedblleft{}On the contrary; by contrast\textquotedblright{}.) Say it once more.
\end{tcolorbox}
